% root = ../paper.tex
\section{$f$-Differential Privacy and Its Basic Properties}
\label{sec:fDP}


%\djs{Replace $\T$ notation.}
%\djs{Revise roadmap after we agree on the content of this section.}



In Section~\ref{sub:trade_off_function}, we give a formal definition of $f$-DP. Section~\ref{sub:gaussian_differential_privacy} introduces Gaussian differential privacy, a special case of $f$-DP. In Section~\ref{sec:conn-with-blackw}, we highlight some appealing properties of this new privacy notation from an information-theoretic perspective. Next, Section~\ref{sub:a_primal_dual_connection_with_} offers a profound connection between $f$-DP and $(\epsilon, \delta)$-DP. Finally, we discuss the group privacy properties of $f$-DP.

% Finally, we discuss connections between $f$-DP and Blackwell's theorem, $(\epsilon, \delta)$-DP, and its ``group privacy'' properties in %Sections~\ref{sec:conn-with-blackw}, \ref{sub:a_primal_dual_connection_with_}, and \ref{sub:group_privacy}, respectively.

% %%Then we discuss its properties and relation with other differential privacy notions such as $(\ep,\delta)$-DP and {\Renyi} DP.
% %%The central novelty of our theory is in the characterization of indistinguishablility: we switch from various divergences to \textit{trade-off functions}, which is
%introduced in \Cref{sub:trade_off_function}. \djs{Here or intro?}
%Before that, let's first clarify a few terms for readers not familiar with differential privacy.
% and then start from the motivation and definition of the central concepts of this paper in \cref{sub:trade_off_function}.

%%As a central idea, differential privacy requires that neighboring datasets should be almost indistinguishable.
%%Such a dataset is often organized as a $n\times d$ matrix, with rows corresponding to individuals. In this view,
%%a neighboring dataset of $\vec{x}$ is a matrix $\vec{x}'$ that differs $\vec{x}$ by one row. For example, $\vec{x}' = (x_1',\ldots, x_n)$ is a neighbor of $\vec{x}$ as %long as $x_1'\neq x_1$. What we just described is replace-one neighboring. There is also add-one/remove-one neighboring.


%%\djs{Changed to dataset. We might want to distinguish between collection and ordered collection. Now we claim a dataset is a collection, but our mathematical %representation is an ordered collection.}
%% \djs{Still not very consistent. I use $X^n$ in Sections 4 and 5 without explicitly saying $X=\R^d$.} Without loss of generality, a dataset $x$ can be thought of as a %collection of $n$ records, each consisting of a $d$-dimensional real-valued vector \ar{Any reason to make it a $d$-dimensional real valued vector? Why not have an %abstract data domain $\mathcal{X}$?}. We write:

Before moving on, we first establish several key pieces of notation from the differential privacy literature.
\begin{itemize}
\item  % (fold)\label{par:dataset}
{\bf Dataset.} A dataset $S$ is a collection of $n$ records, each corresponding to an individual. Formally, we write the dataset as $S = (x_1,\ldots, x_n)$, and an individual $x_i \in X$ for some abstract space $X$. Two datasets $S' = (x'_1,\ldots, x'_n)$ and $S$ are said to be \textit{neighbors} if they differ in exactly one record, that is, there exists an index $j$ such that $x_i = x_i'$ for all $i \ne j$ and $x_j \ne x_j'$.

%\wjs{say $Y$}\djs{Removed ``by appropriate permuting''. We do not explicit say our algorithms are invariant under the action of permutation group, and we don't have to.}

% Replace-$k$ neighboring corresponds to group privacy (see \Cref{sub:group_privacy}). In order to minimize such detail, it is helpful to make the following abstraction:
% % \begin{definition} \label{def:dataset}
% \begin{center}
% 	A \emph{dataset space} is an abstract set $X$ equipped with a symmetric binary relation $\sim$.
% \end{center}
% \djs{Putting it into a definition environment breaks the ``Database'' and ``Mechanism'' paragraph environments.}
% % \end{definition}
%%$X$ is a collection of possible instances of datasets. The neighboring relation is encoded in $\sim$. Note that a transitive $\sim$ typically does not lead to a meaningful privacy notion. \djs{More should be said on this in intro.}
% \ar{In fact, its not transitive in any interesting setting, right?}
% paragraph dataset (end)

\item {\bf Mechanism.}  A mechanism $M$ refers to a randomized algorithm that takes as input a dataset $S$ and releases some (randomized) statistics $M(S)$ of the dataset in some abstract space $Y$. For example, a mechanism can release the average salary of individuals in the dataset plus some random noise.

% (fold) \label{par:mechanism}
%\ar{Is this definition useful/correct? It doesn't have to output ``summary statistics'' -- the range can be anything. All we are saying is mechanism = algorithm. }
% . For some readers it may be helpful to think of a mechanism $M$ as a Markov kernel or conditional probability. When there is no ambiguity, the notation often intended for mappings, $M:X\to Y$, is abused to mean that $M$ is a mechanism/randomized algorithm whose input and output are elements of $X$ and $Y$ respectively.
% More precisely, when $M:X\to Y$ is fed with two neighboring datasets $x\sim x'$, the results $M(x)$ and $M(x')$ are distributions on $Y$. For example, $M:\R^{n\times d}\to\R^d$ may take the average of each column and add Gaussian noise to the average vector. $M(x)$ and $M(x')$ will be two $d$-dimensional Gaussian with different means in this case.
% % paragraph mechanism (end)

\end{itemize}

%%\subsection{Trade-off functions and Testing Regions} % (fold)

\input{fDP/tradeoff}
\input{fDP/GDP}
\input{fDP/post}
\input{fDP/dual}
\input{fDP/group}



%\input{fDP/group}
% subsection a_primal_dual_connection_with_ (end)
%%% Local Variables:
%%% mode: latex
%%% TeX-master: "../paper"
%%% End:
